\section{Ejecución en Tiempo Real}
\label{sec:runtime}

\subsection{El loop de eventos: \texttt{rclpy.spin()}}

Una vez que el nodo está construido y todos sus suscriptores, publicadores y timers registrados, la llamada \texttt{rclpy.spin(node)} entrega el control al loop de eventos de ROS~2. La Terminal~\ref{lst:spin_loop} muestra su comportamiento conceptual:

\begin{figure}[H]
\centering
\begin{lstlisting}[style=bash, caption={Comportamiento conceptual de \texttt{rclpy.spin()}.}, label={lst:spin_loop}]
# Pseudocodigo del loop de eventos
while rclpy.ok():
    esperar_evento()          # bloquea hasta que algo ocurre
    ejecutar_callback()       # procesa el evento
    # repetir
\end{lstlisting}
\end{figure}

El proceso permanece bloqueado en este loop indefinidamente. El código del usuario solo vuelve a ejecutarse cuando algún evento lo activa: la llegada de un mensaje en un tópico suscrito, o el disparo de un timer. Fuera de esos momentos, el proceso duerme y no consume CPU.

\subsection{Suscripciones y callbacks de mensajes}

Una suscripción registra un callback que ROS~2 llama automáticamente cada vez que llega un mensaje en el tópico indicado, como muestra la Terminal~\ref{lst:odom_sub}:

\begin{figure}[H]
\centering
\begin{lstlisting}[style=python, caption={Registro de una suscripción a odometría.}, label={lst:odom_sub}]
self.odom_sub = self.create_subscription(
    Odometry,           # tipo del mensaje
    'odom',             # nombre del topico
    self.odom_callback, # funcion a llamar
    10                  # queue size (QoS History)
)
\end{lstlisting}
\end{figure}

El argumento \texttt{10} es el tamaño de la cola: si los mensajes llegan más rápido de lo que el callback los procesa, se acumulan hasta 10 mensajes antes de empezar a descartarlos. DDS gestiona la entrega desde el publicador (el nodo de odometría de Gazebo) hasta este callback, incluso si están en procesos diferentes.

El callback de suscripción típicamente solo actualiza el estado interno del nodo:

\begin{figure}[H]
\centering
\begin{lstlisting}[style=python, caption={Callback de suscripción: actualiza estado, no computa control.}, label={lst:odom_cb}]
def odom_callback(self, msg: Odometry):
    self.odom = msg   # guarda el mensaje completo
    # NO computar control aqui: llega a frecuencia variable
\end{lstlisting}
\end{figure}

La razón de no computar el control dentro del callback de suscripción es que los mensajes de odometría llegan a la frecuencia que el publicador decida (o que la red permita), no a la frecuencia que el algoritmo de control necesita. Mezclar la lógica de control con la llegada de datos produce un nodo con frecuencia de operación impredecible.

\subsection{Timers: el ciclo de control}

Un timer registra un callback que se ejecuta periódicamente con frecuencia determinista. La Terminal~\ref{lst:ctrl_timer} muestra cómo crear el loop de control a 20~Hz:

\begin{figure}[H]
\centering
\begin{lstlisting}[style=python, caption={Registro de un timer de control a 20~Hz.}, label={lst:ctrl_timer}]
self.timer = self.create_timer(
    0.05,               # periodo en segundos (20 Hz)
    self.control_loop   # funcion a llamar
)
\end{lstlisting}
\end{figure}

El timer garantiza que \texttt{control\_loop} se ejecute cada 50~ms independientemente de cuándo llegaron los últimos mensajes. Esta separación entre la \textit{adquisición de datos} (callbacks de suscripción) y el \textit{cómputo de control} (callback del timer) es el patrón fundamental en nodos de control con ROS~2.

\subsection{El patrón completo: adquisición y control separados}

La Terminal~\ref{lst:control_pattern} muestra el patrón completo de un nodo de control tal como se estructura en este proyecto:

\begin{figure}[H]
\centering
\begin{lstlisting}[style=python, caption={Patrón completo de nodo de control: suscriptores actualizan estado, timer computa y publica.}, label={lst:control_pattern}]
def odom_callback(self, msg):
    self.odom = msg          # actualiza estado: 50-100 Hz

def scan_callback(self, msg):
    self.scan = msg          # actualiza estado: 10 Hz

def control_loop(self):      # ejecuta a 20 Hz exactos
    # guardia: datos aun no disponibles
    if self.odom is None or self.scan is None:
        return

    # computo de control con datos mas recientes
    cmd = self.pva.compute(self.odom, self.scan)

    # publicacion del comando de velocidad
    self.cmd_pub.publish(cmd)
\end{lstlisting}
\end{figure}

Este patrón resuelve el problema de la asincronía inherente a los sistemas distribuidos: el nodo no bloquea esperando datos, sino que usa los datos más recientes disponibles en el momento en que el timer dispara. Si el LiDAR llega a 10~Hz y el control corre a 20~Hz, cada par de ciclos de control usará el mismo scan, lo cual es completamente aceptable.

\subsection{Modelo de concurrencia: SingleThreadedExecutor}

Por defecto, \texttt{rclpy.spin()} utiliza un \texttt{SingleThreadedExecutor}, lo que significa que todos los callbacks (suscriptores y timers) se ejecutan en el mismo hilo, uno a la vez. Esto tiene dos consecuencias:

Primero, el acceso a variables compartidas como \texttt{self.odom} y \texttt{self.scan} es seguro sin necesidad de locks: el executor garantiza que nunca habrá dos callbacks ejecutándose simultáneamente. Segundo, si un callback tarda más que su periodo (por ejemplo, el loop de control tarda 60~ms siendo el periodo 50~ms), el siguiente disparo del timer se retrasa. Los callbacks deben ser lo suficientemente rápidos para no acumular retraso.

Si se necesita paralelismo real (procesamiento de imagen pesado mientras el control sigue corriendo), existe \texttt{MultiThreadedExecutor} con \texttt{ReentrantCallbackGroup}, pero requiere manejo explícito de acceso concurrente a las variables compartidas.

\subsection{El rol de DDS en la entrega de mensajes}

DDS es la capa de transporte subyacente de ROS~2. Cuando un nodo arranca, DDS anuncia su existencia en la red y otros nodos que suscriben al mismo tópico se conectan automáticamente, sin configuración manual de IPs o puertos. Para la transmisión, DDS serializa los mensajes al formato wire, los transmite (localmente vía memoria compartida, o en red vía UDP) y los deserializa en el receptor.

El usuario de ROS~2 no interactúa con DDS directamente; \texttt{rclpy} lo abstrae completamente. Sin embargo, entender que DDS existe explica por qué los nodos se descubren automáticamente y por qué es posible comunicar nodos en máquinas diferentes con solo configurar \texttt{ROS\_DOMAIN\_ID}.
